% Copyright 2026 The terCAD team. All rights reserved.
% Use of this content is governed by a CC BY-NC-ND 4.0 license that can be found in the LICENSE file.

\subsection{PHP}
\markboth{Bootcamping}{PHP} \label{php}

A set of tools for personal web pages initiated the development of PHP in 1994. Over time, it
evolved into a general-purpose language with C-style syntax that supports dynamic and gradual
typing. PHP supports procedural, object-oriented (including inheritance, polymorphism, and
encapsulation), and functional programming styles.

\paragraph{Zend Engine} compiles PHP code into an intermediate representation (opcodes) that is
executed by a virtual machine. OPcache stores the compiled opcodes in shared memory, which
allows it to skip parsing on subsequent runs. The just-in-time (JIT) compiler, introduced in
version 8.0, translates hot paths into native machine code. Memory is managed by reference
counting with a cycle-collecting garbage collector. Traditionally, PHP follows a shared-nothing
model: each request starts in a clean state and releases everything at the end.

\paragraph{Composer} is a dependency manager that controls libraries and their versions via
\q{composer.json}-file (and locks them in \q{composer.lock}) by using public
(\href{https://packagist.org}{https://packagist.org}) and private
(\q{"repositories": [\{"type": "vcs", "url": "https://..."\}]}) sources. It also generates an
autoloader that maps namespaces to directories (PSR-4 standard), so classes are loaded on the
first usage without any manual \q{require}-statement.

\paragraph{Entire End-to-End Stack} PHP is well-equipped for full-stack development. Laravel and
Symfony provide a robust foundation for web applications and APIs. Livewire brings reactive
interfaces without leaving PHP. Swoole, RoadRunner, and FrankenPHP keep the application in
memory as a long-running process. The database layer is covered by PDO (MySQL, PostgreSQL,
SQLite, etc.) and ORMs (Eloquent, Doctrine).


\subsubsection{Checking Primitives} \label{php-prim}

Scalar types are \q{int}, \q{float}, \q{bool}, and \q{string}; compound -- \q{array},
\q{object}, \q{callable}, and \q{iterable}; special -- \q{null} and \q{resource}.
Variables are prefixed by \q{\$}-sign:

\begin{lstlisting}
$pi = 3.141592; // round($pi, 2) => 3.14
$text = 'Some content'; // strlen($text) => 12
$tmp = 1; $tmp = 'test'; // Changeable type
// Double quotes and heredoc interpolate variables
$sample = "{$pi} items";
$long = <<<TXT
    Hello, {$text}!
    TXT;
// Concatenation by '.'
$line = 'some text...' . '... another content' . PHP_EOL;
\end{lstlisting}

\noindent The \q{array} is a universal ordered map that plays roles of a list, a dictionary, a
stack, and a queue at once. For specific cases, the Standard PHP Library (SPL) offers
\q{SplStack}, \q{SplQueue}, \q{SplObjectStorage}, \q{SplFixedArray}, and \q{SplPriorityQueue}:

\newpage
\begin{lstlisting}
$data = [1, 2, 3];
// end($data) returns 3
$map = ['test' => 123];
// array_values($map) returns [123]
$scope = array_unique(['a', 'b', 'a']);
// count($scope) returns 2
// Spread and destructuring
$merged = [...$data, 4];    // [1, 2, 3, 4]
[$first, , $third] = $data; // 1, 3
['test' => $value] = $map;  // 123
// Since 8.4
array_find($data, fn($v) => $v > 1); // 2
\end{lstlisting}


\subsubsection{Observing Types}

By default, PHP converts types implicitly (\emph{"weak" mode}). Declaring
\q{declare(strict\_types=1);} at the top of a file turns it into a strict mode, where passing
a wrong scalar type throws \q{TypeError}. Types can be declared for arguments, returns,
properties, and (\emph{since 8.3}) class constants:

\begin{lstlisting}
declare(strict_types=1);

function p(float $a, ?string $c = null): string
{
    return sprintf('%.2f %s', $a, $c ?? 'EUR');
}
\end{lstlisting}

\noindent Static analyzers (\q{PHPStan}, \q{Psalm}, \q{Mago}) extend native types by
annotations, such as generics:

\begin{lstlisting}
/** @return list<array{name: string, age: int}> */
function users(): array { /* ... */ }
\end{lstlisting}


\subsubsection{Matching Conditions}

The \q{match}-expression (\emph{since 8.0}) is a strict (\q{===}) alternative to \q{switch} that
returns a value and throws \q{UnhandledMatchError} if no condition fits:

\begin{lstlisting}
$view = match ($this->activeTab) {
    'chat', 'forum' => new Communication(),
    'map' => new Map(),
    default => new Home(),
};
// 'match(true)' to evaluate conditions
$label = match (true) {
    $hp <= 0 => 'dead',
    $hp < 30 => 'wounded',
    default => 'healthy',
};
\end{lstlisting}

\noindent Null-handling operators keep checks short:

\begin{lstlisting}
// null if any part is null
$name = $user?->profile?->name;
// default 'Guest' if missing
$nick = $input['nick'] ?? 'Guest';
// assign only if $cache is null
$cache ??= loadCache();
\end{lstlisting}


\subsubsection{Operating with Dates}

\q{DateTimeImmutable} returns a new object on every change (preferred over the mutable
\q{DateTime}), while \q{DateInterval} and \q{DatePeriod} describe durations and ranges. Time
zones are handled natively:

\begin{lstlisting}
$start = new DateTimeImmutable('2026-01-01', new DateTimeZone('Europe/Warsaw'));
$end = $start->modify('+1 month');
$days = $start->diff($end)->days; // 31
echo $end->format('D, M j, Y'); // Sun, Feb 1, 2026
// Parsing without exceptions on failure
$value = DateTimeImmutable::createFromFormat('Y-m-d', 'not-a-date'); // false
// Next, instead of failing, will calculate overflow
// for out-of-range values: 2027-02-14
$overflow = DateTimeImmutable::createFromFormat('Y-m-d', '2026-13-45');
\end{lstlisting}

\noindent Localized formatting is provided by \q{IntlDateFormatter} (\q{intl}-extension), and
\q{Carbon}-library (bundled with Laravel) adds a fluent interface with an ability to freeze the
current time in tests (\q{Carbon::setTestNow()}).


\subsubsection{Declaring Input Arguments}

Arguments can be positional, optional (with defaults), named (\emph{since 8.0}), and variadic:

\begin{lstlisting}
function add(string $value, ?int $id = null, string ...$tags): void {}

add('test');                    // Optional 'id'
add('test', 123, 'a', 'b');     // Variadic 'tags'
add(value: 'test', id: 1);      // Named arguments
add(...['value' => 'test']);    // Unpacking
\end{lstlisting}


\subsubsection{Initializing Classes}

Each class has a single parent (\q{extends}), may implement multiple interfaces
(\q{implements}), and reuse code horizontally via \q{trait}s. Constructor promotion,
\q{readonly} properties, property hooks, and asymmetric visibility (\emph{since 8.4})
reduce boilerplate:

\newpage
\begin{lstlisting}
final class Hero
    extends Character
    implements JsonSerializable
{
    use HasInventory; // trait

    public function __construct(
        // immutable
        public readonly string $name,
        // read publicly, write privately
        public private(set) int $level = 1,
    ) {}

    public string $title {
        get => "{$this->name} [{$this->level}]";
    }

    public function jsonSerialize(): array
    {
        return [
            'name' => $this->name,
            'level' => $this->level
        ];
    }
}

// No extra parentheses since 8.4
echo new Hero('Tom')->title; // Tom [1]
\end{lstlisting}

\noindent Class modifiers: \q{abstract} (prevents instantiation), \q{final} (prohibits
inheritance), and \q{readonly} (all properties are immutable). Visibility is controlled by
\q{public}, \q{protected}, and \q{private}; \q{static} grants access without instantiation.


\subsubsection{Enumerating Classifications}

Enums (\emph{since 8.1}) are classes with a predefined set of cases. Backed enums keep a scalar
value, and both can contain methods and implement interfaces:

\begin{lstlisting}
enum Status: string
{
    case Active = 'active';
    case Pending = 'pending';

    public function label(): string
    {
        return ucfirst($this->value);
    }
}

Status::from('active'); // Status::Active
Status::tryFrom('unknown'); // null (no exception)
Status::cases(); // [Status::Active, Status::Pending]
\end{lstlisting}


\subsubsection{Observing Functional Programming}

Functions are first-class citizens: they can be stored in variables, passed as arguments, and
returned. Closures capture outer variables via \q{use}, while arrow functions (\q{fn}) capture
them automatically:

\begin{lstlisting}
$prefix = 'INFO';
$logger = function (string $msg) use ($prefix) { echo "$prefix: $msg"; };
$double = fn(int $x) => $x * 2;

// First-class callable syntax (since 8.1)
$upper = strtoupper(...);
array_map($double, [1, 2, 3]); // [2, 4, 6]

// Pipe operator (since 8.5)
$result = '  hello ' |> trim(...) |> $upper; // HELLO
\end{lstlisting}

\noindent Generators (\q{yield}) produce values lazily, which is handy to iterate over large
files or responses:

\begin{lstlisting}
function lines(string $path): Generator
{
    $file = fopen($path, 'r');
    while (($line = fgets($file)) !== false) {
        yield $line;
    }
    fclose($file);
}
\end{lstlisting}


\subsubsection{Operating Asynchronously}

PHP executes code synchronously, but \q{Fiber}s (\emph{since 8.1}) allow to pause and resume a
function, which is the foundation of the \q{Revolt} event loop, \q{Amp} (v3), and the
\q{react/async}-package of \q{ReactPHP}. Parallel HTTP requests are also available
through \q{curl\_multi} (or \q{Http::pool} in Laravel), and heavy tasks are commonly moved into
queues:

\begin{lstlisting}
$fiber = new Fiber(function (string $x): string {
    $y = Fiber::suspend($x . ' paused');
    return $y . ' resumed';
});
echo $fiber->start('Task'); // Task paused
$fiber->resume('Task');
echo $fiber->getReturn(); // Task resumed
\end{lstlisting}


\subsubsection{Overloading Behavior}

"Magic methods" declare a custom behavior for built-in operations:

\begin{itemize}
  \item \q{\_\_toString} -- string representation of an object;
  \item \q{\_\_invoke} -- allows an object to be called as a function;
  \item \q{\_\_get}, \q{\_\_set}, \q{\_\_isset}, \q{\_\_unset} -- access to inaccessible
  properties;
  \item \q{\_\_call}, \q{\_\_callStatic} -- calls of undefined methods
  (widely used by Laravel Facades);
  \item \q{\_\_serialize}, \q{\_\_unserialize}, \q{\_\_clone} -- control serialization and
  cloning.
\end{itemize}

\begin{lstlisting}
class Model
{
    public function __call(string $v, array $a): mixed
    {
        if (str_starts_with($v, 'save')) {
            return "Saving with: " . implode(',',$a);
        }
        throw new BadMethodCallException($v);
    }
}
(new Model)->saveIssue('#1'); // Saving with: #1
\end{lstlisting}

\noindent Operators themselves cannot be overloaded in userland code (only by extensions, as
\q{GMP} or \q{BCMath}).


\subsubsection{Using Attributes}

Attributes (\emph{since 8.0}) attach structured metadata to classes, methods, properties, and
arguments; they are read via the Reflection API and power routing, validation, and testing in
frameworks:

\begin{lstlisting}
#[Attribute(Attribute::TARGET_METHOD)]
final class Route
{
    public function __construct(public string $path) {}
}

abstract class Controller
{
    public function index(): void {}
}

class ChatController extends Controller
{
    #[Route('/chat')]
    #[\Override] // since 8.3: fatal error if the parent has no such method
    public function index(): void {}
}

$attrs = (new ReflectionMethod(ChatController::class, 'index'))->getAttributes(Route::class);
echo $attrs[0]->newInstance()->path; // /chat
\end{lstlisting}


\subsubsection{Managing Components}

Namespaces isolate names, and the \q{use}-statement imports them with aliases to resolve
collisions. Composer's autoloader resolves a class to its file based on the namespace:

\begin{lstlisting}
namespace App\Services;

use App\Models\Notification;
use Illuminate\Support\Facades\Http;
use Carbon\Carbon as Date;   // alias
use function Laravel\Prompts\text;
\end{lstlisting}

\noindent The mapping is configured in \q{composer.json}:

\begin{lstlisting}[language=java]
"autoload": {
    "psr-4": {
       "App\\":"app/",
       "Database\\Factories\\":"database/factories/",
       "Database\\Seeders\\":"database/seeders/"
    }
}
\end{lstlisting}


\subsubsection{Asserting in Tests}

\q{PHPUnit} and \q{Pest} (built on top of PHPUnit) are the most common testing frameworks:

\begin{lstlisting}
// strict comparison (===)
$this->assertSame(0, $result);
// loose comparison (==)
$this->assertEquals([1, 2], $result);

$this->assertStringContainsString('text', $result);
$this->assertMatchesRegularExpression('/^\w+$/', $result);
$this->assertContains('a', $list);
$this->assertArrayHasKey('id', $map);
$this->assertCount(3, $list);
$this->assertGreaterThan(0, $result);
$this->assertNull($result);
// assertTrue, assertEmpty, ...
$this->assertInstanceOf(Hero::class, $result);
$this->expectException(TypeError::class);
// Pest-style
expect($result)->toBe(0)->not->toBeNull();
\end{lstlisting}

\noindent In addition, PHPUnit mocks (\q{createMock}) and \q{Mockery} declare expectations to
ensure that some other methods were called (or not) during the test execution:

\begin{lstlisting}
// PHPUnit: replace a dependency and verify its usage
$proxy = $this->createMock(AppProxyProvider::class);
$proxy->expects($this->once())
    ->method('boot')
    ->with($this->stringContains('/ch/chmain.php'))
    ->willReturn('var sesid = "abc";');
$proxy->expects($this->never())->method('decode');

$html = $proxy->boot('https://example.com/ch/chmain.php');
$this->assertSame('abc', ListenStream::sessionIdFrom($html));

// Mockery: the same expectations in a fluent way
$proxy = Mockery::mock(AppProxyProvider::class);
$proxy->shouldReceive('boot')
    ->once()
    ->with('/ch/chmain.php')
    ->andReturn('<html></html>');
$proxy->shouldNotReceive('decode');
\end{lstlisting}
